%% =====================================================================
%%  T-05-01  The Real Objection
%%  -------------------------------------------------------------------
%%  One idea per file. Core is mandatory. Slots in canonical order.
%%  Max 700 words. No layout commands. No filler. New source material
%%  for this entry goes in sources/<BOOK>/T-05-01.tex -- never by
%%  rewriting this file (docs/RULES.md R0.1).
%% =====================================================================
%% @kind: topic
%% @id: T-05-01
%% @title: The Real Objection
%% @subtitle: The stated reason is a position, not a fact
%% @chapter: c05
%% @order: 10
%% @status: stable
%% @sources: B1
%% @updated: 2026-09-19

\topic{T-05-01}{The Real Objection}{The stated reason is a position, not a fact}

\begin{core}
People give reasons that are presentable. The obstacle actually blocking them is usually something else --- money, fear, a third party's opinion, an earlier commitment --- and until it is identified, every attempt to move them is aimed at something that is not in the way.
\end{core}
\begin{mechanism}
The stated reason is chosen partly to be defensible and partly because it is the reason the person can articulate. Real obstacles are frequently unarticulable: an insecurity, an obligation to someone who is not in the room, a private estimate of their own capacity. Answering the stated reason therefore produces a strange pattern --- each objection, once fully answered, is replaced by a new one --- and that pattern is itself the diagnostic. A person whose real objection has been met stops generating objections; a person whose stated objection has been met produces another.
\end{mechanism}
\begin{tells}
  \item Objections arrive serially: each resolved one is followed by a fresh one.
  \item The stated reason would not, if true, prevent the action.
  \item The reason changes depending on who is present.
  \item They defend the reason more vigorously than they want the outcome.
  \item Relief rather than engagement when you stop pressing.
\end{tells}
\begin{moves}
  \item Stop answering objections and start collecting them. Three in a row is a pattern; one is noise.
  \item Ask what would have to be true for them to want this. The answer is usually nearer the real obstacle than anything they have said so far.
  \item Offer a face-saving route to the truth: make it respectable for the real reason to be money, fear or someone else's opinion.
  \item Address the real obstacle directly once identified, and let the stated reasons fall away on their own.
  \item If the real obstacle cannot be addressed, stop. Persuading past it produces compliance that collapses later.
\end{moves}
\begin{counters}
  \item Notice your own serial objections. If you have given three reasons for refusing something, the fourth is being manufactured to cover the first.
  \item Answer the real one, not the presented one, even when doing so is uncomfortable.
  \item Where someone is mining your objection, give the true reason once and hold it. A single stable refusal is harder to work on than a shifting one.
  \item Ask for the decision to be deferred. Obstacle-mining depends on the conversation continuing.
\end{counters}
\begin{cost}
Searching for a hidden objection can become a refusal to accept no. A stated reason is sometimes the reason, and the serial-objection test is the only thing separating legitimate persistence from pressure --- once the answer is a settled refusal rather than a sequence of pretexts, continuing is the coercion described in T-19-01 and not persuasion. The diagnostic also fails on the literal counterpart from T-03-01, who has no hidden layer to find.
\end{cost}

\sources{B1}
\seesources
\seealso{T-05-02, T-07-02, T-16-01, T-24-02}
